\section{Experimental Design Details}
\label{sec:appendix_design}

\subsection{Constant Valuations}

Each bidder has a constant valuation $v_i=1.0$. Table~\ref{tab:exp1params} details the 10 factorial factors and fixed parameters. The $2^{10-1}$ Resolution~V half-fraction yields 512 cells, each replicated twice for 1{,}024 observations. The bid grid resolution is fixed at 11 discrete actions; a discretisation sensitivity analysis across grid sizes of 6, 11, and 21 is reported in Section~\ref{sec:appendix_robustness}.

\begin{table}[h]
\centering
\caption{Parameter settings for Experiment~1 (Constant Valuations).}
\label{tab:exp1params}
\begin{tabular}{l l l}
\toprule
\textbf{Factor} & \textbf{Low ($-1$)} & \textbf{High ($+1$)}\\
\midrule
Auction type & Second-price & First-price \\
Learning rate ($\alpha$) & 0.01 & 0.10 \\
Discount factor ($\gamma$) & 0.0 & 0.95 \\
Reserve price ($r$) & 0.0 & 0.5 \\
Initialisation & Zeros & Optimistic \\
Exploration & $\varepsilon$-greedy & Boltzmann \\
Update mode & Synchronous & Asynchronous \\
Number of bidders ($n$) & 2 & 4 \\
Information feedback & None (stateless) & Previous winning bid \\
Decay type & Linear & Exponential \\
\midrule
\multicolumn{3}{l}{\textbf{Fixed parameters}} \\
\midrule
Bid grid resolution & \multicolumn{2}{l}{11 discrete actions\footnotemark} \\
Number of episodes & \multicolumn{2}{l}{100{,}000} \\
\bottomrule
\end{tabular}
\footnotetext{Bid grid resolution is held constant at 11 levels. A discretisation sensitivity analysis varying the grid size is reported in Section~\ref{sec:appendix_robustness}.}
\end{table}

\subsection{Affiliated Valuations with Q-Learning}

Bidders receive private signals $s_i\in[0,1]$, and their valuations become interdependent through an affiliation parameter $\eta\in[0,1]$ via $v_i=(1-0.5\,\eta)\,s_i + 0.5\,\eta\,\bigl(\tfrac{1}{n-1}\sum_{j\neq i}s_j\bigr)$. This lets $\eta=0$ represent purely private values and $\eta=1$ represent strong common-value elements. Table~\ref{tab:exp2params} highlights the key parameters.

\begin{table}[h]
\centering
\caption{Parameter settings for Experiment~2 (Affiliated Valuations + Q-learning).}
\label{tab:exp2params}
\begin{tabular}{l l l}
\toprule
\textbf{Factor} & \textbf{Low ($-1$)} & \textbf{High ($+1$)}\\
\midrule
Auction type & Second-price & First-price \\
Number of bidders ($n$) & 2 & 4 \\
State information & Signal only & Signal + winner \\
\midrule
\textbf{Factor} & \multicolumn{2}{l}{\textbf{Levels}} \\
\midrule
Affiliation ($\eta$) & \multicolumn{2}{l}{0, 0.5, 1 (linear + quadratic contrasts)} \\
\midrule
\multicolumn{3}{l}{\textbf{Fixed parameters (from Experiment~1 results)}} \\
\midrule
Learning rate ($\alpha$) & \multicolumn{2}{l}{0.1} \\
Discount factor ($\gamma$) & \multicolumn{2}{l}{0.95} \\
Exploration & \multicolumn{2}{l}{$\varepsilon$-greedy} \\
Update mode & \multicolumn{2}{l}{Asynchronous} \\
Decay type & \multicolumn{2}{l}{Linear} \\
Reserve price ($r$) & \multicolumn{2}{l}{0.0} \\
Initialisation & \multicolumn{2}{l}{Zeros} \\
Bid grid resolution & \multicolumn{2}{l}{11 discrete actions} \\
Number of episodes & \multicolumn{2}{l}{100{,}000} \\
\bottomrule
\end{tabular}
\end{table}

\subsection{Affiliated Valuations with Bandits}

This experiment preserves the affiliated valuation model from Experiment\,2 but replaces Q-learning with contextual bandit methods. Table~\ref{tab:exp3params} summarises the parameters.

\begin{table}[H]
\centering
\caption{Parameter settings for Experiment~3 (Affiliated Valuations + Bandits).}
\label{tab:exp3params}
\begin{tabular}{l l l}
\toprule
\textbf{Factor} & \textbf{Low ($-1$)} & \textbf{High ($+1$)}\\
\midrule
Algorithm & LinUCB & CTS \\
Auction type & Second-price & First-price \\
Number of bidders ($n$) & 2 & 4 \\
Reserve price ($r$) & 0.0 & 0.3 \\
Exploration intensity & Low & High \\
Context richness & Signal only & Signal + winner \\
Regularisation ($\lambda$) & 0.1 & 5.0 \\
\midrule
\textbf{Factor} & \multicolumn{2}{l}{\textbf{Levels}} \\
\midrule
Affiliation ($\eta$) & \multicolumn{2}{l}{0, 0.5, 1 (linear + quadratic contrasts)} \\
\midrule
\multicolumn{3}{l}{\textbf{Fixed parameters}} \\
\midrule
Bid grid resolution & \multicolumn{2}{l}{11 discrete actions} \\
Number of rounds & \multicolumn{2}{l}{100{,}000} \\
\bottomrule
\end{tabular}
\end{table}

\subsection{Budget-Constrained Pacing}

Each of $n \in \{2, 4\}$ advertisers participates in $D = 100$ episodes, each comprising $T = 1{,}000$ single-item auctions; budgets regenerate between episodes while dual variables persist (warm-starting). Valuations follow a log-normal model with bidder-specific asymmetry. All bidders use multiplicative dual pacing, computing bids as $v_t / \mu_t$ (value-maximisers) or $v_t / (1 + \mu_t)$ (utility-maximisers), subject to a hard budget cap. Budgets are set as $B_i = m \cdot \mathbb{E}[v_i] \cdot T$ where $m$ is a budget multiplier that varies as an experimental factor. Table~\ref{tab:exp4aparams} details the factor levels.

\begin{table}[H]
\centering
\caption{Parameter settings for Experiment~4a (Budget-Constrained Pacing).}
\label{tab:exp4aparams}
\begin{tabular}{l l l}
\toprule
\textbf{Factor} & \textbf{Low ($-1$)} & \textbf{High ($+1$)}\\
\midrule
Auction type & Second-price & First-price \\
Bidder objective & Value-maximizer & Utility-maximizer \\
Number of bidders ($n$) & 2 & 4 \\
Budget multiplier ($m$) & 0.25 (tight) & 1.0 (loose) \\
Reserve price ($r$) & 0.0 & 0.3 \\
Value dispersion ($\sigma$) & 0.1 (low) & 0.5 (high) \\
\midrule
\multicolumn{3}{l}{\textbf{Fixed parameters}} \\
\midrule
Episodes ($D$) & \multicolumn{2}{l}{100} \\
Rounds per episode ($T$) & \multicolumn{2}{l}{1{,}000} \\
Log-normal $\mu$ range & \multicolumn{2}{l}{$[0.5, 1.5]$} \\
Initial $\mu_0$ & \multicolumn{2}{l}{1.0} \\
Burn-in episodes & \multicolumn{2}{l}{10} \\
\bottomrule
\end{tabular}
\end{table}

\subsection{PI Controller Pacing}

Each of $n \in \{2, 4\}$ advertisers participates in $D = 100$ episodes, each comprising $T = 1{,}000$ single-item auctions; budgets regenerate between episodes while the PI multiplier $\lambda$ persists (warm-starting). All bidders use a PI pacing controller that updates $\lambda$ additively based on cumulative spending error (Section~\ref{sec:appendix_algorithms}, Equations~\ref{eq:pid_error}--\ref{eq:pid_update}). The aggressiveness parameter $a$ jointly scales both the proportional gain $K_P = 0.30a$ and integral gain $K_I = 0.05a$. Table~\ref{tab:exp4bparams} details the factor levels.

\subsection{Factor Level Calibration}
\label{sec:factor_calibration}

Factorial designs code each factor at two levels ($-1$ and $+1$) and attribute outcome differences to factor identity. This subsection documents calibration analyses for factors where asymmetries between levels arise and explains why the factorial estimates remain interpretable for the research questions posed. When the low-to-high contrast produces different effective changes across levels, the estimated main effect conflates factor identity with intensity; the analyses below quantify this for each affected factor.

\subsubsection{Exploration Strategy Calibration (Experiment~I)}

Experiment~I compares $\varepsilon$-greedy and Boltzmann exploration using identical parameter schedules: both decay from $1.0$ to $0.01$ over the first 90\% of training. However, these raw parameters control qualitatively different mechanisms. For $\varepsilon$-greedy, $\varepsilon$ directly sets the probability of a uniform random action, producing Shannon entropy $H = (1-\varepsilon)\cdot 0 + \varepsilon \cdot \ln(K) = \varepsilon \ln K$ for $K$ actions when Q-values are well-separated. For Boltzmann, the temperature $\tau$ scales the Q-value differences before the softmax, and the resulting entropy depends on both $\tau$ and the Q-value spread. At identical raw parameter values, Boltzmann can maintain substantially higher or lower entropy than $\varepsilon$-greedy depending on the Q-value magnitudes.

To address this, Boltzmann exploration uses Q-range normalized temperatures. At each decision point, the effective temperature is
\[
  \tau_{\text{eff}} = \tau \cdot \max\bigl(\Delta Q,\; \epsilon\bigr), \quad \Delta Q = \max_a Q(s,a) - \min_a Q(s,a),
\]
where $\tau$ follows the same decay schedule as $\varepsilon$ and $\epsilon = 10^{-8}$ prevents division-related issues when Q-values are uniform (e.g., at initialisation). This normalization ensures the softmax probabilities $\pi(a|s) \propto \exp(Q(s,a)/\tau_{\text{eff}})$ depend on the relative spacing of Q-values rather than their absolute scale, making the Boltzmann temperature comparable in effect to $\varepsilon$-greedy across different discount factors and training stages. When Q-values are well-separated, $\Delta Q$ is large and the effective temperature scales proportionally, preventing premature exploitation; when Q-values converge, $\Delta Q$ shrinks and the agent naturally sharpens its policy.

\subsubsection{Decay Schedule Asymmetry (Experiment~I)}

The decay type factor contrasts linear decay ($\varepsilon_t = 1 - t/t_{\max}$) with exponential decay ($\varepsilon_t = \varepsilon_0 \cdot (\varepsilon_{\min}/\varepsilon_0)^{t/t_{\max}}$). At the midpoint of training, exponential decay retains a higher exploration rate than linear decay (ratio approximately 1.6:1), producing more total integrated exploration over the training horizon. This asymmetry is by design: the factor tests whether front-loaded exploration (linear) or sustained exploration (exponential) better supports learning. The factorial model's main effect for decay type captures this intended contrast, and the small estimated effect sizes across response variables confirm that decay shape is a secondary consideration relative to factors such as auction format and discount factor.

\subsubsection{Exploration Intensity Asymmetry (Experiment~III)}

In Experiment~III, the exploration intensity factor maps to algorithm-specific parameters: for LinUCB, the confidence bound width $c \in \{0.5, 2.0\}$ (a 4$\times$ ratio), and for Contextual Thompson Sampling, the posterior variance $\sigma^2 \in \{0.1, 1.0\}$ (a 10$\times$ ratio). These ratios reflect the different units and scales of the two exploration mechanisms. LinUCB explores by inflating the estimated reward with $c\sqrt{\mathbf{x}^\top A^{-1}\mathbf{x}}$, where the uncertainty term shrinks with data; a 4$\times$ increase in $c$ roughly quadruples the exploration bonus. Thompson Sampling explores through posterior sampling variance $\sigma^2 A^{-1}$; a 10$\times$ increase broadens the sampling distribution proportionally. The factorial model captures this asymmetry through the algorithm $\times$ exploration intensity interaction term, which estimates how the marginal effect of increasing exploration differs between the two algorithms. A separate entropy measurement using empirical action frequencies over 1{,}000-round windows confirms that the effective exploration intensity change differs between algorithms, as expected from the parameter ratio asymmetry.

\subsubsection{Memory Decay (Experiment~III)}

The memory decay factor $\delta$ controls how much weight historical observations receive relative to recent ones. At $\delta = 1.0$ (no decay), both LinUCB and Thompson Sampling accumulate all past data, and after $T$ rounds the effective sample size is $T$. At $\delta < 1$, the effective sample size plateaus at approximately $1/(1-\delta)$. This factor is motivated by \citet{Douglas2024}, who show that deterministic convergence in bandit learners drives collusive outcomes, and by \citet{Russac2019}, who provide regret guarantees for discounted linear bandits in non-stationary environments. Specific levels and calibration rationale are presented after the exploratory simulations described in Section~\ref{sec:factor_calibration}.

\subsubsection{Structural Factor Asymmetries}

Several factors across experiments involve inherently asymmetric level contrasts that reflect the experimental questions rather than calibration failures.

In Experiment~I, the information feedback factor changes the state space from 1 state (no feedback) to $K = 11$ states (discretised previous winning bid), an 11-fold increase in Q-table size. This is the intended contrast: the factor tests whether additional market information aids or hinders convergence. Similarly, the Q-table initialisation factor (zeros vs.\ optimistic at $1/(1-\gamma)$) produces level-dependent ranges that interact with the discount factor, ranging from 1.0 at $\gamma = 0$ to 20.0 at $\gamma = 0.95$. This interaction is captured by the init $\times$ gamma term in the factorial model.

In Experiment~II, the state information factor similarly expands the state space from $K$ signal bins to $K^2$ signal-winner pairs, a contrast designed to test the curse of dimensionality in Q-learning.

In Experiment~IVa, the bidder objective factor contrasts value-maximisation ($b = v/\mu$) with utility-maximisation ($b = v/(1+\mu)$), two structurally different bidding rules drawn from the price of anarchy literature. The asymmetry is definitional, not parametric.

In Experiment~IVb, the aggressiveness factor scales the PI controller gains by a factor of $10\times$ (from $a = 0.3$ to $a = 3.0$). The proportional gain $K_P$ ranges from $0.09$ to $0.90$ and the integral gain $K_I$ from $0.015$ to $0.15$. Higher aggressiveness produces faster budget consumption and more reactive pacing, while lower aggressiveness yields conservative, slow-adjusting behaviour. The $10\times$ ratio ensures meaningful separation in controller dynamics without pushing either level into degeneracy (budget exhaustion in the first few rounds or near-zero spending).

In Experiment~III, the regularisation parameter $\lambda$ takes values $0.1$ and $5.0$, a 50-fold ratio. Because $\lambda$ controls the initial precision matrix $A = \lambda I$, the inverse $A^{-1}$ that enters both LinUCB's confidence bound ($c\sqrt{\mathbf{x}^\top A^{-1}\mathbf{x}}$) and Thompson Sampling's posterior covariance ($\sigma^2 A^{-1}$) is 50 times larger at $\lambda=0.1$ than at $\lambda=5.0$. Through the square root in LinUCB's bonus, this translates to approximately a $\sqrt{50}\approx 7$-fold difference in initial exploration intensity. The nonlinearity diminishes as data accumulates and the outer-product updates dominate $A$, but it is strongest during early training when exploration matters most. The LGBM robustness comparison reported in Section~\ref{sec:appendix_robustness} validates that the linear OLS model remains an adequate approximation despite this nonlinearity; the $R^2$ gap between the linear and nonparametric models is small for all response variables.

In all cases, the factorial design's main effects and interactions correctly decompose the outcome variation attributable to each factor and their combinations. The asymmetries described above are features of the experimental questions, not confounds; they are properly captured by the orthogonal contrast structure of the $2^k$ design.

\subsection{Summary Statistics and Response Variables}
\label{sec:appendix_summary}

\subsubsection{Experiment~1}

\input{tables/exp1_summary}

\subsubsection{Experiment~2}

Table~\ref{tab:exp2_responses} defines the additional outcome metrics computed in Experiment~2 beyond the shared response variables defined in the main paper. These metrics exploit the affiliated valuation structure and known Bayesian Nash Equilibrium benchmarks to decompose revenue shortfalls and characterise bidding behaviour. Signal slope ratio compares observed bidding to the theoretical linear equilibrium, while winner's curse frequency measures overbidding directly.

\begin{table}[H]
\centering
\caption{Additional response variables for Experiment~2.}
\label{tab:exp2_responses}
\begin{tabular}{l l}
\toprule
\textbf{Metric} & \textbf{Definition} \\
\midrule
Bid-to-value median & Median of payment / winner valuation in final window \\
Winner's curse frequency & Fraction of wins where payment exceeds winner's valuation \\
Bid dispersion & Mean within-round standard deviation of bids \\
Signal slope ratio & OLS slope of bid on signal, divided by BNE slope \\
\bottomrule
\end{tabular}
\end{table}

\input{tables/exp2_summary}

\subsubsection{Experiment~3}

\input{tables/exp3_summary}

\subsubsection{Experiment~4b}

Experiment~4b uses the same response variables as Experiment~4a (Table~\ref{tab:exp4a_responses}), with the dual variable CV computed from $\lambda$ history. The bid suppression ratio uses a competitive benchmark of $\lambda = 1.0$ (full-value bidding) for both auction formats. All remaining metrics are defined identically to Experiment~4a.

\input{tables/exp4b_summary}
